\chapter{Introduction générale}
\label{ch:intro}

\section{Contexte et motivation}

Le métier d'avocat au Maroc suppose une production écrite dense : mémoires, conclusions et surtout \textbf{plaidoiries} destinées à l'audience. Dans de nombreux cabinets, le circuit reste artisanal : notes manuscrites, dictée informelle, ressaisie par un collaborateur, mise en page Word tardive. Ce circuit est lent, source d'erreurs de copie, et peu traçable.

Le Maroc compte plusieurs milliers d'avocats répartis sur dix-neuf barreaux. Le besoin n'est pas un « chatbot juridique » grand public : c'est un \textbf{outil de cabinet} qui accélère la transformation de la parole de l'avocat en document exploitable, tout en respectant la déontologie\footnote{L'avocat demeure seul responsable du contenu déposé ou plaidé. LexHub est un instrument d'aide à la rédaction, non un substitut au conseil juridique.}.

C'est dans ce cadre que le stage s'est déroulé chez \textbf{WEBVOICE} (Tanger), sous tutelle de l'\textbf{ENSA Tétouan} (Université Abdelmalek Essaâdi). Encadrement professionnel : M.~Hatim Benhaddou. Tuteur pédagogique : M.~Mohammed Al Achhab.

Le produit livré s'appelle \textbf{LexHub} (\url{https://lexhub.space}). Le sujet --- une plateforme d'assistance IA pour avocats --- se concrétise par un atelier de \textbf{plaidoirie vocale} au sein d'un espace cabinet sécurisé, complété par dossiers, documents, planning, activité et équipe.

\section{L'entreprise d'accueil --- WebVoice}
\label{sec:webvoice}

\textbf{WebVoice} (\url{https://webvoice.pro}) est une société de Tanger spécialisée dans la \textbf{voix sur IP (VoIP)}, le développement \textbf{web} et \textbf{mobile}, ainsi que les sites e-commerce. Son positionnement --- téléphonie IP, applications métier, accompagnement de bout en bout --- explique naturellement l'accueil d'un projet LegalTech centré sur la \textbf{voix de l'avocat} et sa transformation en document.

Selon sa présentation publique, l'entreprise structure ses prestations en quatre étapes : \textbf{Analyse} des besoins, \textbf{Concept}, \textbf{Mise en œuvre}, puis \textbf{Suivi}. Cette méthode a structuré le stage : cadrage du cœur plaidoirie, conception de l'espace cabinet, implémentation incrémentale, déploiement et consolidation.

LexHub s'inscrit donc dans le prolongement de l'expertise WebVoice (voix, web, SaaS) appliquée au cabinet d'avocats marocain bilingue.

\section{Problématique}

Comment concevoir et déployer une plateforme SaaS \textbf{multi-cabinets} permettant à un avocat marocain de dicter sa plaidoirie en arabe, d'obtenir un texte structuré et un export officiel, de collaborer avec un stagiaire \emph{sans} lui exposer les brouillons non validés, le tout sans que l'IA n'invente le fond juridique ?

\section{Objectifs}

\begin{enumerate}[label=\textbf{O\arabic*.}]
  \item Réaliser le pipeline \textbf{voix $\rightarrow$ texte $\rightarrow$ document} (cinq étapes guidées dans l'interface).
  \item Offrir un espace cabinet : dossiers (domaine + filtre), documents, planning lié aux affaires, équipe, activité.
  \item Isoler strictement les données par cabinet (\textit{multi-tenant}\footnote{Chaque cabinet est un locataire : ses dossiers, fichiers et plaidoiries sont invisibles des autres cabinets.}).
  \item Distinguer les rôles : avocat, stagiaire / assistant, avocat admin, administrateur plateforme WebVoice.
  \item Permettre l'\textbf{assignation} d'une plaidoirie à un assistant, l'épinglage et la validation explicite avant partage.
  \item Mettre en production une version utilisable (TLS, e-mails, temps réel, export Word/PDF) sur \url{https://lexhub.space}.
\end{enumerate}

\section{Méthodologie}

Travail en cycles courts : cadrage (cahier des charges, maquettes), implémentation incrémentale, tests par rôle, déploiement Docker sur VPS, retours des encadrants (premier suivi : recentrage sur voix~$\rightarrow$~plaidoirie, robustesse STT/LLM, benchmarking OK/Non~OK). Le présent rapport décrit \textbf{uniquement la version actuellement déployée}.

\section{Plan du rapport}

\begin{itemize}[leftmargin=1.4em]
  \item Le chapitre~\ref{ch:etat-art} présente l'étude comparative (benchmarking) des solutions existantes.
  \item Le chapitre~\ref{ch:analyse} formalise acteurs, cas d'utilisation UML et exigences.
  \item Le chapitre~\ref{ch:conception} détaille architecture, pipelines et diagrammes UML (cas d'utilisation, séquences, classes, activité, déploiement).
  \item Le chapitre~\ref{ch:realisation} décrit la réalisation, le guide d'utilisation et la sécurité.
  \item Les chapitres~\ref{ch:tests} et~\ref{ch:conclusion} traitent des tests, du déploiement, de l'état CI/CD, des risques et des perspectives d'évolution.
\end{itemize}
